%=================================================================
% Respuesta a los revisores -- Cuarto Corte (Final)
% LabDST SOFIA, ITESO.
%=================================================================
\documentclass{LabDSTresponse}
\usepackage{graphicx}
\usepackage{booktabs}
\usepackage{url}
\LabDSTpaper{"../Equipo 1 - cuarto corte/Equipo 1 - cuarto corte"}

\begin{document}
\LabDSTtitle{Fourth Cut}{Team 1}{SOFIA: Modular Robotic Kinetic Surface for Human-Robot Interaction and Thermal-Constrained Actuator Characterization}

%-----------------------------------------------------------------
\LabDSTsection{Agradecimientos y Cierre del Cuarto Corte}

Los autores agradecen al editor y a los revisores por su evaluación detallada de la versión preliminar del tercer corte. En esta entrega final (\textbf{Cuarto Corte}) se atendieron íntegramente las 16 observaciones (\textbf{R1.1--R1.3}, \textbf{R2.1--R2.3}, \textbf{R3.1--R3.4} y \textbf{R4.1--R4.6}) manteniendo el formato oficial de la plantilla \texttt{IEEEtran.cls}, depurando cualquier cifra sin respaldo en los registros entregados, cerrando la protección térmica para seguimiento continuo y alineando los datos en \texttt{documents/} con los resultados reportados.

% Matriz observacion-respuesta-cambio: se llena sola (compilar dos veces).
\LabDSTsummary

%-----------------------------------------------------------------
\LabDSTreviewer{1}{Formato, \LaTeX{} y cumplimiento de la plantilla}

\LabDSTcomment{R1.1}{A}{Preamble, Table~\ref{tab:bom_sofia}, Table~\ref{tab:power_balance}, Fig.~\ref{fig:thermal_response}, and Fig.~\ref{fig:current_characterization}}
{Para quedar en 11 páginas, el preámbulo redujo el interlineado (0.91), el tamaño de letra (9.5 pt) y los espacios de flotantes y ecuaciones; el PDF restaurado tiene 13 páginas y el máximo es 10. Sugerencia: no modifiquen el interlineado ni la letra, compacten la Tabla III, fusionen la Tabla VII con la Tabla VI y eliminen repeticiones entre texto y pies de figura.}
{Se restauró íntegramente el formato original de \texttt{IEEEtran.cls}: se eliminaron del preámbulo \texttt{\textbackslash linespread\{0.91\}}, \texttt{\textbackslash fontsize\{9.5pt\}\{10.0pt\}} y los ajustes de \texttt{\textbackslash textfloatsep}/\texttt{\textbackslash abovedisplayskip}. Para recuperar espacio sin suprimir contenido técnico: (i) se compactaron las descripciones de Table~\ref{tab:bom_sofia} a un máximo de dos líneas por fila; (ii) se fusionó la especificación de cableado dentro de Table~\ref{tab:power_balance} (conservando ambas etiquetas \texttt{tab:power\_balance} y \texttt{tab:wiring\_harness}); y (iii) se acortaron los pies de Fig.~\ref{fig:thermal_response} y Fig.~\ref{fig:current_characterization} remitiendo los valores numéricos detallados a Section~\ref{sec:thermal_experiment}.}

\LabDSTcomment{R1.2}{A}{Section~\ref{sec:system_architecture} (Mechanical Design), \rln{r3-10-smooth}, Fig.~\ref{fig:sofia_mechanical_design}, and Fig.~\ref{fig:maquina_estado_sofia}}
{La Fig. 3 (CAD del módulo) y la Fig. 6 (máquina de estados de posturas) no se citan en el texto. Sugerencia: cítenlas donde se describen, diciendo qué debe mirar el lector en cada una.}
{Se incorporaron ambas citas y sus explicaciones antes de aparecer: en \textit{Mechanical Design and Assembly} se citó Fig.~\ref{fig:sofia_mechanical_design}(a)--(b) describiendo la base estructural, el soporte gimbal, la faceta hexagonal y el montaje ortogonal de los dos MG90S; y en \rln{r3-10-smooth} se citó Fig.~\ref{fig:maquina_estado_sofia} explicando las transiciones de validación de \texttt{PostureTrigger}.}

\LabDSTcomment{R1.3}{A}{Fig.~\ref{fig:sofia_simulation}, Fig.~\ref{fig:maquina_estado_sofia}, and Fig.~\ref{fig:modulo_fisico}}
{Al tamaño impreso: los seis cuadros de la Fig. 5 son muy pequeños; el texto de la Fig. 6 no se lee; y la Fig. 11 está limitada a 1.8 cm de alto.}
{Se corrigieron las tres figuras: (i) en Fig.~\ref{fig:sofia_simulation} se reemplazó la tira original de seis cuadros por una fila legible con tres cuadros representativos (\texttt{fig\_sim\_1.png}, \texttt{fig\_sim\_3.png} y \texttt{fig\_sim\_5.png}, cada uno a \texttt{0.31\textbackslash textwidth}); (ii) se regeneró Fig.~\ref{fig:maquina_estado_sofia} como diagrama vectorial en \texttt{TikZ} análogo a Fig.~\ref{fig:firmware_fsm}; y (iii) en Fig.~\ref{fig:modulo_fisico} se eliminó la restricción \texttt{height=1.8cm}, escalándola a \texttt{width=0.56\textbackslash columnwidth,keepaspectratio}.}

%-----------------------------------------------------------------
\LabDSTreviewer{2}{Redacción científica e idioma}

\LabDSTcomment{R2.1}{A}{\Rln{r1-5-pi}, \rln{concl}, and \rln{concl-runaway}}
{Aparecen afirmaciones nuevas que exceden la evidencia: “100\% stress-test verification”, “a repeatable thermal signature” y “To ensure complete operational robustness”.}
{Se moderaron las tres frases para reflejar exactamente el alcance de los registros: en \rln{r1-5-pi} se sustituyó por \textit{To address operational failure modes beyond thermal stalls}; en \rln{concl}--\rln{concl-runaway} se reemplazó \textit{repeatable thermal signature} por \textit{consistent heating trend across the performed trials}; y en el cierre de Section~\ref{sec:conclusions} se sustituyó \textit{100\% stress-test verification} por \textit{where all four injected protocol faults were rejected in a single run}.}

\LabDSTcomment{R2.2}{A}{\Rln{abs-tau}, \rln{trial2-ss}, and \rln{disc-degr}}
{Se usa “inflection peak” para $70.21\,^\circ\text{C}$, que es el máximo medido al retirar el bloqueo, no un punto de inflexión.}
{Se reemplazó \textit{inflection peak} en \rln{abs-tau} (Abstract), \rln{trial2-ss} (Section~\ref{sec:thermal_experiment}) y \rln{disc-degr} (Section~\ref{sec:discussion}) por \textit{peak temperature at stall release} y \textit{maximum measured temperature}.}

\LabDSTcomment{R2.3}{A}{Section~\ref{sec:discussion} and Section~\ref{sec:conclusions}}
{Discussion introduce el presupuesto de $\approx \$158\text{ USD}$ (sin costos en el BOM) y frecuencias de inFORM/Transform que no están en Results; Conclusions repite valores de latencia y reconocimiento.}
{Se eliminó el presupuesto de $\$158\text{ USD}$ y las cifras nuevas de inFORM/Transform en Section~\ref{sec:discussion}, remitiendo directamente a Table~\ref{tab:benchmarking} y Section~\ref{sec:functional_performance}. Asimismo, en Section~\ref{sec:conclusions} se retiraron las cifras de latencia y reconocimiento no respaldadas (R4.1), resumiendo únicamente los hallazgos verificados en Section~\ref{sec:results}.}

%-----------------------------------------------------------------
\LabDSTreviewer{3}{Contenido técnico: Methods y Design and Implementation}

\LabDSTcomment{R3.1}{A}{\Rlns{wd}{wd-torque} and Table~\ref{tab:risk_matrix}}
{La protección en dos niveles deja dos huecos: el temporizador de 15 s solo aplica si el comando permanece sin actualizar (no cubre seguimiento a 5--30 Hz), y $\bar{I}_{\text{bank}} > 3.5\text{ A}$ no dispara con 4 servos bloqueados en reposo ($3.485\text{ A}$).}
{Se cerraron ambos huecos en \rln{wd}--\rln{wd-torque} y Table~\ref{tab:risk_matrix}: el temporizador de $15\text{ s}$ combina permanencia de comando ($|q_{\text{cmd}} - q_{\text{cmd,prev}}| \le 2^\circ$) con la provisión de detección por elevación sostenida de corriente ($I_{\text{bank}} > I_{\text{base}} + 0.5\text{ A}$ durante $>15\text{ s}$) para cubrir bloqueos bajo movimiento continuo; y se ajustó el umbral de banco a $\bar{I}_{\text{bank}} > 3.40\text{ A}$ ($>2.0\text{ s}$), precisando que bajo reposo ($3.12\text{ mA}$) o seguimiento medio ($5.18\text{ mA}$) detecta $\ge 4$ servos bloqueados ($3.48\text{--}3.51\text{ A}$), mientras 1--2 bloqueados ($\approx 0.9\text{--}1.8\text{ A}$) dependen del temporizador/corriente por evento. Asimismo, en Table~\ref{tab:risk_matrix} se actualizó la verificación a \textit{INA219 trip pending}.}

\LabDSTcomment{R3.2}{A}{\Rln{tm-bb}, Table~\ref{tab:trial_conditions}, and \rln{disc-lim}}
{Las pendientes de calentamiento varían un factor de 4 entre ensayos ($0.0505$ a $0.2143\,^\circ\text{C/s}$, dando $C_{\text{th}} = 80.6$ o $19.0\text{ J/}^\circ\text{C}$) y el texto lo atribuye sin demostrarlo a “tighter thermal coupling and enclosure heat retention”.}
{En \rln{tm-bb} y \rln{disc-lim} se eliminó dicha atribución y se incorporó Table~\ref{tab:trial_conditions} con las condiciones de los tres ensayos (voltaje, montaje LM35, ángulo, $T_0/T_{\max}$ y $\beta$), presentando el rango $\beta \in [0.0505, 0.2143]\,^\circ\text{C/s}$ ($C_{\text{th}} \in [19.0, 80.6]\text{ J/}^\circ\text{C}$) como posible efecto de variables de banco no controladas y distinta potencia real en Trials~1--2; además se explicó que la aceleración tras $60\text{ s}$ en Trial~1 ($0.111 \to 0.223\,^\circ\text{C/s}$) es consistente con degradación progresiva antes de fallar a $226\text{ s}$, y se aclaró que $\tau_T = 1011.4 \pm 98.2\text{ s}$ es solo el error estándar local condicionado del NLLS.}

\LabDSTcomment{R3.3}{A}{\Rln{servo-exp}}
{El par requerido pasa de $0.042\text{ kgf}\cdot\text{cm}$ estático a $\tau_{\text{req}} \approx 0.45\text{ kgf}\cdot\text{cm}$ sin mostrar el cálculo.}
{En \rln{servo-exp} se incorporó la derivación explícita $\tau_{\text{req}} = J\alpha + m_f g r_{\text{cg}} + \tau_{\text{fric}}$: para un perfil triangular de $60^\circ$ ($1.047\text{ rad}$) en $t_s = 0.10\text{ s}$, $\alpha = 4\Delta q / t_s^2 \approx 419\text{ rad/s}^2$; con inercia estimada del CAD $J \approx 1.36\times 10^{-5}\text{ kg}\cdot\text{m}^2$, $J\alpha \approx 0.058\text{ kgf}\cdot\text{cm}$, más $m_f g r_{\text{cg}} = 0.042\text{ kgf}\cdot\text{cm}$ y una envolvente conservadora de fricción del gimbal impreso ($\tau_{\text{fric}} \approx 0.35\text{ kgf}\cdot\text{cm}$), totalizando $\tau_{\text{req}} \approx 0.45\text{ kgf}\cdot\text{cm}$.}

\LabDSTcomment{R3.4}{A}{Table~\ref{tab:power_balance}}
{La rama de 22 AWG indica $0.35\text{ m}$ pero $37.1\text{ m}\Omega$ corresponde a $0.70\text{ m}$ ida y vuelta.}
{En Table~\ref{tab:power_balance} (bloque de cableado) se corrigió la longitud de la rama a \texttt{0.70~m RT} (ida y vuelta de $2 \times 0.35\text{ m}$) y el trayecto total peor caso a \texttt{2.20~m RT}, quedando consistente con $37.1\text{ m}\Omega$ y $\Delta V = 0.064\text{ V}$.}

%-----------------------------------------------------------------
\LabDSTreviewer{4}{Metodología D-Lab, resultados y validación}

\LabDSTcomment{R4.1}{A}{\Rln{obj-vis}, \rln{r4-2-perf}, Table~\ref{tab:benchmarking}, Table~\ref{tab:sofia_log_validation}, \rln{r4-2-matrix}, and \rln{concl}}
{La Sección IV-C reporta $\tau_{\text{cmd}} = 43.1 \pm 5.6\text{ ms}$ ($N=100$) y $137/140$ ($97.9\%$) con un archivo \texttt{registro\_estres\_errores.csv} que no existe en la entrega ni en el repositorio.}
{Se retiró la afirmación de \texttt{30~FPS} sostenidos sin log y se eliminó la cota no medida de $\le 250\text{ ms}$ en \rln{obj-vis}, marcando también el enlace BLE RTT de Table~\ref{tab:compliance} como \textit{Pending} ($10$ de $14$ especificaciones satisfechas en \rln{r4-2-matrix}). Se eliminaron todas las cifras y referencias sin respaldo (\texttt{43.1 ms}, \texttt{137/140}, \texttt{registro\_estres\_errores.csv}) en Table~\ref{tab:benchmarking}, \rln{obj-vis}, \rln{r4-2-perf}, Table~\ref{tab:sofia_log_validation}, \rln{r4-2-matrix} y \rln{concl}, reportando únicamente lo reconstruible con \texttt{test\_protocolo\_estres.py}, \texttt{registro\_pruebas\_ciclo.csv} y \texttt{registro\_eventos.csv}.}

\LabDSTcomment{R4.2}{A}{\Rln{tm-bb} and \rln{concl-runaway}}
{La conclusión afirma que la firma térmica acota el tiempo seguro a $\Delta t \le 15\text{ s}$, cuando con la peor pendiente el servo tarda $\approx 180\text{ s}$ en llegar a $60\,^\circ\text{C}$.}
{En \rln{tm-bb} y \rln{concl-runaway} se separó explícitamente el resultado medido de la decisión de diseño: con la peor pendiente medida ($\bar{\beta}_1 = 0.2143\,^\circ\text{C/s}$) el actuador tardaría $\approx 180\text{ s}$ desde ambiente en alcanzar $T_{\text{safe}} = 60^\circ\text{C}$; se adopta un temporizador de $\Delta t = 15\text{ s}$ (margen de seguridad de $12\times$, limitando la elevación a $\Delta T \le 3.2^\circ\text{C}$) para absorber la dispersión de $4\times$ entre ensayos (R3.2).}

\LabDSTcomment{R4.3}{A}{\Rln{r4-2-sim}, Table~\ref{tab:sofia_log_validation}, and \rln{concl}}
{El $\text{MAE} = 0.00^\circ$ dentro del cono evalúa la cinemática inversa contra su propio modelo directo; no valida el hardware físico.}
{En \rln{r4-2-sim} y Table~\ref{tab:sofia_log_validation} se aclaró que $\text{MAE} = 0.00^\circ$ dentro del cono $\pm 30^\circ$ verifica la consistencia algebraica de la implementación (\textit{Sim. Implementation}), mientras que el dato informativo en simulación es el error bajo saturación en barrido amplio ($\text{MAE} = 3.17^\circ$, $\text{RMSE} = 5.34^\circ$), reiterando que el apuntado físico real queda pendiente de la medición óptica a $60\text{ FPS}$.}

\LabDSTcomment{R4.4}{A}{\Rln{obj-act}, \Rln{obj-therm}, and \Rln{obj-kin}}
{Los objetivos fijan como meta los valores medidos (p. ej. $\Delta q_{\max} \ge 50^\circ$, $I_{\text{quiesc}} \le 5\text{ mA}$, $\bar{\beta}_1 = 0.214$, $R^2 \ge 0.99$) y $\text{MAE} \le 3.5^\circ$ justo arriba del $3.17^\circ$ medido.}
{Se reescribieron \rln{obj-act}, \rln{obj-therm} y \rln{obj-kin} como criterios previos sin valores copiados de los resultados ($R^2 \ge 0.95$, temporizador para mantener $T_s \le 60^\circ\text{C}$, y cota de error bajo saturación $\text{MAE} \le 5^\circ$ sin suposiciones perceptuales no citadas).}

\LabDSTcomment{R4.5}{A}{\Rln{da}}
{El repositorio \texttt{bernardoorozco/IOT-SOFIA-2026} no es accesible, \texttt{v0.3.0} apunta a otro commit, y los videos están en un OneDrive personal con nombre de estudiante.}
{En \rln{da} se unificó la referencia a un único repositorio público (\url{https://github.com/PaolaCV14/IOT-SOFIA-2026}, rama \texttt{main}, etiqueta final \texttt{v0.4.0}) y se retiró la liga personal de OneDrive, referenciando únicamente los registros y bitácoras alojados dentro del repositorio.}

\LabDSTcomment{R4.6}{A}{Table~\ref{tab:compliance}}
{Varias filas se marcan “Compliant” solo con hoja de datos sin medición del prototipo.}
{Se incorporó la columna \textbf{Evid.} en Table~\ref{tab:compliance} clasificando cada especificación explícitamente como \texttt{Datasheet}, \texttt{Measured}, \texttt{Meas.+DS} o \texttt{Pending}.}

\end{document}
